%%=============================================================================
%% Inleiding
%%=============================================================================

\chapter{\IfLanguageName{dutch}{Inleiding}{Introduction}}%
\label{ch:inleiding}

Streamingdiensten zijn in de huidige samenleving nauwelijks nog weg te denken. Platformen zoals \emph{Netflix}, \emph{Disney+} en \emph{Spotify} bieden audiovisuele content aan door middel van abonnementsformules, waardoor gebruikers deze op elk moment en op diverse toestellen kunnen bekijken of beluisteren. Ook in België wordt er steeds meer gestreamd \autocite{BNA2025}, en daar speelt de Vlaamse entertainmentindustrie actief op in: de voorbije jaren lanceerden verschillende Vlaamse zenders en productiehuizen hun eigen platformen, waaronder \emph{Streamz}, \emph{VRT MAX} en \emph{VTM GO}, waardoor het aanbod aanzienlijk groter werd. De groei van het aantal streamingdiensten en de dalende populariteit van live televisie \autocite{Digimeter2024} leiden ertoe dat steeds meer Vlamingen gebruikmaken van al dan niet betalende streamingdiensten \autocite{BNA2025}.

\vspace{1em}
Maar deze groeiende populariteit brengt nieuwe gevaren met zich mee. Uit onderzoek van \textcite{FODEconomie2024} blijkt dat veel streamingdiensten zogenaamde dark patterns toepassen. In de literatuur worden deze patronen ook vaak aangeduid als \emph{Dark Patterns}. Volgens deze studie gaat het om misleidende ontwerptechnieken die online gebruikers subtiel beïnvloeden om bepaalde keuzes te maken die zij anders niet zouden maken. Voorbeelden hiervan zijn onder meer moeilijk vindbare annulatieknoppen of vooraf aangevinkte selectievakjes \autocite{FODEconomie2024}.

%De inleiding moet de lezer net genoeg informatie verschaffen om het onderwerp te begrijpen en in te zien waarom de onderzoeksvraag de moeite waard is om te onderzoeken. In de inleiding ga je literatuurverwijzingen beperken, zodat de tekst vlot leesbaar blijft. Je kan de inleiding verder onderverdelen in secties als dit de tekst verduidelijkt. Zaken die aan bod kunnen komen in de inleiding~\autocite{Pollefliet2011}:
%
%\begin{itemize}
%  \item context, achtergrond
%  \item afbakenen van het onderwerp
%  \item verantwoording van het onderwerp, methodologie
%  \item probleemstelling
%  \item onderzoeksdoelstelling
%  \item onderzoeksvraag
%  \item \ldots
%\end{itemize}

\vspace{1em}

\section{\IfLanguageName{dutch}{Probleemstelling}{Problem Statement}}%
\label{sec:probleemstelling}

\vspace{1em}

Hoewel bedrijven al lange tijd het koopgedrag van consumenten proberen te beïnvloeden door middel van reclame, behouden consumenten hun keuzevrijheid. Aan de hand van verkregen productinformatie, demonstraties of prijsvergelijkingen kunnen zij zelfstandig beslissen om al dan niet een aankoop te doen. Bovendien zijn er strenge juridische regels die reclame reguleren om consumenten te beschermen \autocite{UNIZO2025}. Reclame mag volgens \textcite{UNIZO2025} onder andere niet misleidend zijn en moet duidelijk herkend worden als reclame. 

\vspace{1em}

Dit vormt precies het spanningsveld bij het gebruik van dark patterns op streamingdiensten: deze zijn ontworpen om de consumenten te misleiden en worden vaak niet herkend als commerciële beïnvloeding \autocite{FODEconomie2024}. Uit onderzoek van \textcite{BongardBlanchy2021} blijkt dat gebruikers van streamingdiensten zoals Netflix vaak moeite hebben met het identificeren van dergelijke patronen. Verder volgt uit deze studie dat zelfs wanneer consumenten zich bewust zijn van het bestaan van dark patterns, hun keuzes nog steeds beïnvloed worden, wat wijst op de subtiele en manipulatieve aard van deze technieken. Daarnaast kunnen deze patronen leiden tot onbedoelde consequenties voor gebruikers, zoals het onbewust afsluiten van een abonnement of het betalen van een verborgen extra kost \autocite{FODEconomie2024}.

\vspace{1em}

Ondanks dat er al veel onderzoek is verricht naar dark patterns vanuit het perspectief van degenen die deze toepassen \autocite{Kitkowska2023,Kollmer2022,OECD2022}, zijn er nog relatief weinig studies die de effecten van deze patronen op gebruikersgedrag en gebruikerservaring onderzoeken \autocite{BongardBlanchy2021}. Daarnaast benadrukt de studie van \textcite{DiGeronimo2020} de nood aan verder onderzoek om bewustwording bij gebruikers te stimuleren, bijvoorbeeld door de ontwikkeling van geautomatiseerde en educatieve tools. 

\vspace{1em}

Om de doelgroep van dit onderzoek te bepalen werd er gekeken naar de leeftijdsgroep die het meest actief gebruikmaakt van streamingdiensten in België. Uit de statistieken van \textcite{StatistiekVlaanderen2025} blijkt dat dit vooral jongvolwassenen zijn tussen 18 en 34 jaar. Hierdoor ligt de focus voor dit onderzoek op Vlaamse gebruikers tussen 18 en 34 jaar die actief gebruikmaken van streamingdiensten en al dan niet interesse tonen in het afsluiten van een abonnement.

\vspace{1em}

\section{\IfLanguageName{dutch}{Onderzoeksvraag}{Research question}}%
\label{sec:onderzoeksvraag}

\subsection{Hoofdonderzoeksvraag}
\label{subsec:hoofdonderzoeksvraag}

\vspace{1em}

Uit de voorgaande bevindingen werd in het kader van dit onderzoek onderstaande onderzoeksvraag opgesteld: \emph{“Hoe kunnen commerciële dark patterns van in Vlaanderen vaak gebruikte streamingdiensten automatisch opgespoord en aan de gebruikers gesignaleerd worden?”}.

\vspace{1em}

\subsection{Deelvragen}
\label{subsec:deelvragen}

\vspace{1em}

Aanvullend werden enkele deelvragen geformuleerd die ondersteuning bieden bij het beter begrijpen van het probleem en het zoeken naar een mogelijke oplossing.

\subsubsection*{Deelvragen probleemdomein}
\begin{enumerate}
    \item Welke dark patterns komen het meest voor op websites en apps van populaire streamingdiensten die door Vlamingen tussen 18 en 34 jaar gebruikt worden?
    \item In welke mate zijn gebruikers zich bewust van de aanwezigheid van dark patterns bij het gebruik van streamingdiensten?
    \item Welke verbanden bestaan er tussen de aanwezigheid van specifieke dark patterns en de businessdoelstellingen van streamingdiensten?
    \item Op welk platform (website of mobiele app) komen dark patterns het meest frequent voor?
    \item Kunnen technologische tools helpen bij het verhogen van bewustzijn rond bepaalde onderwerpen?
\end{enumerate}

\subsubsection*{Deelvragen oplossingsdomein}
\begin{enumerate}
    \item Welke wet- en regelgeving, richtlijnen of aanbevelingen bestaan er rond het gebruik van dark patterns in België en de Europese Unie?
    \item Welke bestaande tools die gebruikers waarschuwen voor dark patterns bestaan er al, en wat zijn hun sterke en zwakke punten?
    \item Op welke manier kunnen dark patterns aan gebruikers gesignaleerd worden zonder dat hun gebruikerservaring negatief beïnvloed wordt?
    \item Welke digitale of interactieve oplossing is geschikt voor het automatisch detecteren en signaleren van dark patterns op streamingdiensten, en welke (niet-)functionele vereisten zijn hiervoor essentieel?
    \item In welke mate voldoet de uitgewerkte digitale oplossing aan de opgestelde requirements en welke verbeterpunten kunnen worden geïdentificeerd?
\end{enumerate}

\vspace{1em}

\section{\IfLanguageName{dutch}{Onderzoeksdoelstelling}{Research objective}}%
\label{sec:onderzoeksdoelstelling}

\vspace{1em}

Als uiteindelijk doel van dit onderzoek wordt er achterhaald welke manier het beste is om dark patterns op streamingplatformen automatisch op te sporen en hoe deze op een gebruiksvriendelijke manier aan gebruikers gesignaleerd kunnen worden. Naast de scriptie wordt als concreet eindresultaat een Proof of Concept (\acrshort{poc}) uitgewerkt, gebaseerd op de bevindingen uit dit onderzoek. Deze \acrshort{poc} wordt tot slot bij een kleine groep jongvolwassenen (18-34 jaar) getest om na te gaan hoe gebruikers interactie ervaren en hoe effectief de signalering van dark patterns wordt opgevat. De resultaten van deze kleine test geven een eerste indicatie over de gebruiksvriendelijkheid van de voorgestelde oplossing.

\vspace{1em}

\section{\IfLanguageName{dutch}{Opzet van deze bachelorproef}{Structure of this bachelor thesis}}%
\label{sec:opzet-bachelorproef}

\vspace{1em}

% Het is gebruikelijk aan het einde van de inleiding een overzicht te
% geven van de opbouw van de rest van de tekst. Deze sectie bevat al een aanzet
% die je kan aanvullen/aanpassen in functie van je eigen tekst.

Om de onderzoeksvraag systematisch te beantwoorden, is deze bachelorproef verdeeld in de volgende hoofdstukken:

\vspace{1em}

In Hoofdstuk~\ref{ch:stand-van-zaken} wordt een overzicht gegeven van de stand van zaken binnen het onderzoeksdomein, op basis van een literatuurstudie. Verder wordt een overzicht opgesteld van bestaande digitale oplossingen die gebruikers kunnen helpen om dark patterns op streamingdiensten te herkennen.

\vspace{1em}

In Hoofdstuk~\ref{ch:methodologie} wordt de methodologie toegelicht en worden de gebruikte onderzoekstechnieken besproken om een antwoord te kunnen formuleren op de onderzoeksvragen.

\vspace{1em}

In Hoofdstuk~\ref{ch:requirementsanalyse} worden de functionele en niet-functionele requirements van de tool gedefinieerd en geprioriteerd op basis van het literatuuronderzoek en interviews, met als doel het bepalen van de criteria voor het \acrshort{mvp}.

\vspace{1em}

In Hoofdstuk~\ref{ch:shortlist} worden de alternatieven uit de longlist geëvalueerd aan de hand van de opgestelde requirements, wat leidt tot een selectie van de meest beloftevolle oplossing.

\vspace{1em}

In Hoofdstuk~\ref{ch:poc} wordt de geselecteerde oplossing uit Hoofdstuk~\ref{ch:shortlist} uitgewerkt tot een eerste \acrshort{mvp}, inclusief ontwerp, implementatie en validatie via gebruikerstesten. Hiervoor wordt een bestaande open-sourcetool aangepast en uitgebreid voor de specifieke use case van dit onderzoek.

\vspace{1em}

In Hoofdstuk~\ref{ch:conclusie}, tenslotte, wordt de conclusie gegeven en een antwoord geformuleerd op de onderzoeksvragen. Daarbij wordt ook een aanzet gegeven voor toekomstig onderzoek binnen dit domein.